\scanpage{21}
Let $\alpha_0\in K_1$, I say that \unsure{we have|we have}
\[
\widehat{\Sigma}_{\alpha_0}\cap\widehat{\Sigma}_i^+=\varnothing
\qquad\text{i.e}\qquad
\Sigma_{\alpha_0}\cap\Sigma_i=\varnothing
\]
(which will complete the proof, for we shall have
\[
U=X\setminus\bigcup_{\alpha\in K}\widehat{\Sigma}_\alpha^+
\qquad\left\{\begin{aligned}
&K=K_1\cup\{i\}\\
&\deleted{\unsure{\Sigma_\alpha}}
\end{aligned}\right.
\]
as in the lemma.

Now \unsure{as} we have
\[
E_i\subset U_{J,i}=X\setminus\mathop{\rlap{\raisebox{.65ex}{$\scriptstyle-$}}\bigcup}\limits_{\alpha\in K_0}\widehat{\Sigma}_\alpha^+
\]
this implies that
\[
E_i\subset X\setminus\widehat{\Sigma}_{\alpha_0}^+=\Sigma_{\unsure{\alpha_0|\alpha}}^-
\]
% margin note written sideways
\margin{\unsure{j:e.|j.e.|i.e.}\diagram{\begin{tikzcd}[ampersand replacement=\&,column sep=0pt,row sep=0.9ex,execute at end picture={
\draw ([xshift=-4.0mm]dTop.west) .. controls +(0.7mm,-1mm) and +(1mm,2mm) .. ([xshift=-4.1mm]dMid.west) .. controls +(-1mm,-.8mm) and +(1mm,.5mm) .. ([xshift=-5mm]dMid.west) .. controls +(1mm,-.5mm) and +(1mm,.8mm) .. ([xshift=-4.1mm,yshift=-1mm]dMid.west) .. controls +(.5mm,-2mm) and +(-.6mm,1mm) .. ([xshift=-4.6mm]dEnd.west);
\draw ([xshift=-.9mm]dMid.west) .. controls +(-.4mm,-1mm) and +(.9mm,.7mm) .. ([xshift=-1.1mm]dE.west) .. controls +(-1mm,-.4mm) and +(.7mm,.1mm) .. ([xshift=-1.8mm]dE.west) .. controls +(.7mm,-.1mm) and +(.8mm,.5mm) .. ([xshift=-1.1mm,yshift=-.5mm]dE.west) .. controls +(.4mm,-1mm) and +(-.5mm,.6mm) .. ([xshift=-.8mm]dEnd.west);
\node[rotate=90,inner sep=0pt] at ([xshift=-2.7mm,yshift=-1mm]dE.west) {$\scriptscriptstyle\Sigma_i^-\,\unsure{\rotatebox{-90}{$!$}|\text{or}|-}\,\Sigma_i^+$};
}]
|[alias=dTop,anchor=west]| \scriptstyle\Sigma_i^+\ \text{or}\ \Sigma_i^-\\
|[anchor=west]| \scriptstyle E_i\\
|[alias=dMid,anchor=west]| \scriptstyle\Sigma_{i\alpha}\\
|[alias=dE,anchor=west]| \scriptstyle E_\alpha\\
|[alias=dEnd,anchor=west]| \scriptstyle\Sigma_{\alpha_0}^+
\end{tikzcd}}}
on the other hand $\alpha_0\in K_1$ \unsure{mean|means}
\[
\Sigma_{\alpha_0}^+\not\subset\Sigma_i^+
\]
which \deleted{\unsure{impli|implies}} also means that $\Sigma_i^+\cap\Sigma_\alpha^+=\varnothing$, q.e.d.

\uline{Corollary of the lemma} For a class

\noindent\begin{tabular}{@{}|p{.96\linewidth}@{}}
$V$ of $\Pi_I^{\unsure{\circ|0}}$, its expression
\[
V=X\setminus\bigcup_{\alpha\in K}\widehat{\Sigma}_\alpha^+
\]
of the preceding lemma is \uline{unique} -- we have
{\small
\[
\left\{\begin{aligned}
&K=\{\alpha\in I\mid\overline V\cap E_\alpha\ne\varnothing\quad\text{i.e.}\quad\deleted{V}E_\alpha\cap\dot V\ne\varnothing\}\\
&\Sigma_\alpha^+\ \text{is the unique component of }X\setminus\Sigma_\alpha\ \text{contained in }X\setminus V
\end{aligned}\right.
\]}
\end{tabular}
% archive sheet number: 20
